\chapter{Como Explicar Intelig{\^e}ncia Artificial aos Estudantes}
\label{cap:explicar-estudantes}

\begin{objetivos}
\begin{itemize}
  \item Estabelecer o contrato pedag{\'o}gico de transpar{\^e}ncia no primeiro dia de aula;
  \item Desconstruir os mitos do or{\'a}culo infal{\'i}vel e da ferramenta in{\'u}til junto aos estudantes;
  \item Traduzir os cinco n{\'i}veis da Escala AIAS para linguagem acess{\'i}vel \`a turma;
  \item Implantar o protocolo de declara{\c c}{\~a}o de uso de IA nas atividades avaliativas.
\end{itemize}
\end{objetivos}

\section{O Contrato Pedag{\'o}gico da Aula Inaugural}

O maior obst{\'a}culo enfrentado por estudantes em rela{\c c}{\~a}o {\`a} intelig{\^e}ncia artificial n{\~a}o {\'e} a complexidade t{\'e}cnica das ferramentas, mas a \textbf{ambiguidade normativa}. Quando a institui{\c c}{\~a}o ou o docente n{\~a}o explicitam com clareza o que {\'e} permitido e o que {\'e} vedado, cria-se um ambiente de inseguran{\c c}a, medo e desonestidade velada \cite{perkins2024aias}.

O primeiro dia de aula {\'e} o momento determinante para estabelecer o \textbf{contrato pedag{\'o}gico de transpar{\^e}ncia}. Em vez de iniciar o semestre amea{\c c}ando puni{\c c}{\~o}es com base em detectores ineficazes, o professor deve assumir a lideran{\c c}a do di{\'a}logo:

\begin{conceitochave}
\textbf{Princ{\'i}pio da Pactua{\c c}{\~a}o Aberta}: A intelig{\^e}ncia artificial faz parte da realidade profissional e acad{\^e}mica contempor{\^a}nea. O objetivo da universidade n{\~a}o {\'e} fingir que a IA n{\~a}o existe, mas ensinar o estudante a utiliz{\'a}-la com discernimento cr{\'i}tico, honestidade intelectual, rigor metodol{\'o}gico e responsabilidade {\'e}tica.
\end{conceitochave}

\section{Desmistificando a IA em Sala de Aula}

Para que os estudantes adotem uma postura madura diante das ferramentas generativas, {\'e} necess{\'a}rio desconstruir dois mitos opostos igualmente danosos: o mito do \emph{or{\'a}culo infal{\'i}vel} (crer cegamente em tudo o que a m{\'a}quina responde) e o mito da \emph{ferramenta in{\'u}til} (desprezar a tecnologia por preconceito).

Uma analogia pedag{\'o}gica amplamente recomendada (trabalhada na skill \path{ia-educacao-pensamento-critico}) {\'e} apresentar o modelo de linguagem como um \textbf{``estagi{\'a}rio hiperveloz, mas extremamente desatento''}:
\begin{itemize}
  \item Ele leu milh{\~o}es de livros e escreve com velocidade espantosa;
  \item Ele n{\~a}o tem compreens{\~a}o real do mundo e n{\~a}o possui senso moral;
  \item Quando n{\~a}o sabe uma resposta, ele prefere inventar uma mentira convincente a admitir que n{\~a}o sabe;
  \item Portanto, entregar o relat{\'o}rio do estagi{\'a}rio ao cliente sem revis{\~a}o minuciosa {\'e} neglig{\^e}ncia profissional grave. O valor do estudante est{\'a} justamente na sua capacidade de auditar, corrigir e direcionar o trabalho.
\end{itemize}

\subsection{Din{\^a}mica Pr{\'a}tica: A ``Ca{\c c}a ao Erro''}

Uma estrat{\'e}gia did{\'a}tica de alto impacto consiste em projetar na primeira semana de aula uma resposta gerada por IA sobre um t{\'o}pico b{\'a}sico da disciplina que contenha uma alucina{\c c}{\~a}o sutil proposital (uma data equivocada, um julgado citado com tese que n{\~a}o existe, um artigo de lei citado com dispositivo incorreto, um teorema com premissas omitidas ou um c{\'o}digo com vazamento de mem{\'o}ria). 

Desafiar a turma a identificar as falhas em pequenos grupos demonstra imediatamente aos discentes que o uso passivo da IA resulta em mediocridade, enquanto o conhecimento conceitual profundo {\'e} indispens{\'a}vel para validar qualquer sa{\'i}da algor{\'i}tmica.

\section{Traduzindo a Escala AIAS para os Estudantes}

A Escala AIAS n{\~a}o deve permanecer como um documento t{\'e}cnico restrito aos professores. Ela deve ser explicada em linguagem direta e afirmativa aos estudantes, conforme sugerido na Tabela~\ref{tab:aias-discente}.

Antes de apresentar a tabela, informe duas coisas separadamente: se o uso de IA generativa {\'e} permitido ou proibido e se seu uso {\'e} opcional ou obrigat{\'o}rio. Tamb{\'e}m indique as tecnologias assistivas e adapta{\c c}{\~o}es autorizadas, que n{\~a}o devem ser confundidas com o uso de IA generativa.

\begin{table}[htbp]
  \centering
  \small
  \caption{Como Explicar os N{\'i}veis AIAS para a Turma}
  \label{tab:aias-discente}
  \begin{tabularx}{\linewidth}{clX}
    \toprule
    \textbf{N{\'i}vel} & \textbf{Nome} & \textbf{Mensagem Pedag{\'o}gica ao Estudante} \\
    \midrule
    \textbf{1} & Sem IA & ``Nesta tarefa, quero avaliar seu racioc{\'i}nio e suas decis{\~o}es sem uso de IA generativa. Tecnologias assistivas e adapta{\c c}{\~o}es previamente autorizadas continuam permitidas.'' \\
    \addlinespace
    \textbf{2} & Planejamento Assistido por IA & ``Voc{\^e} pode usar a IA para debater ideias, montar esquemas e sugerir organiza{\c c}{\~o}es de t{\'o}picos. Contudo, o texto final e os c{\'o}digos devem ser redigidos integralmente por voc{\^e}.'' \\
    \addlinespace
    \textbf{3} & Colabora{\c c}{\~a}o com IA & ``Voc{\^e} pode usar a IA para sugerir corre{\c c}{\~o}es e refinamentos em trechos espec{\'i}ficos do seu trabalho, desde que documente explicitamente quais partes contaram com aux{\'i}lio da ferramenta.'' \\
    \addlinespace
    \textbf{4} & IA Integral & ``Nesta atividade, o enunciado exige o uso da IA ao longo do processo. Voc{\^e} ser{\'a} avaliado pela qualidade das decis{\~o}es, pela auditoria das respostas, pela documenta{\c c}{\~a}o do processo e pelo dom{\'i}nio cr{\'i}tico do resultado final.'' \\
    \addlinespace
    \textbf{5} & Explora{\c c}{\~a}o de IA & ``Vamos explorar solu{\c c}{\~o}es in{\'e}ditas e interdisciplinares. Voc{\^e} dirigir{\'a} o processo, registrar{\'a} as contribui{\c c}{\~o}es da IA e continuar{\'a} respons{\'a}vel pela valida{\c c}{\~a}o e pelo resultado final.'' \\
    \bottomrule
  \end{tabularx}
\end{table}

\section{O Protocolo de Transpar{\^e}ncia e Atribui{\c c}{\~a}o}

Para atividades nos n{\'i}veis AIAS~2, 3, 4 ou 5, o plano de ensino deve instituir uma cl{\'a}usula simples e transparente de \textbf{Declara{\c c}{\~a}o de Uso de IA}:

\begin{tcolorbox}[colback=white, colframe=boxborder, arc=1mm, title={Modelo de Declara{\c c}{\~a}o de Uso de IA a ser Anexada pelo Estudante}]
\small
\textbf{Declara{\c c}{\~a}o de Transpar{\^e}ncia Acad{\^e}mica}:\\
``Neste trabalho (N{\'i}vel AIAS \underline{\hspace{1cm}}), utilizei a ferramenta \underline{\hspace{3cm}} (vers{\~a}o \underline{\hspace{1.5cm}}) com a seguinte finalidade pedag{\'o}gica: \underline{\hspace{6cm}}.\\
Os principais \emph{prompts} utilizados foram arquivados e coloco-me {\`a} disposi{\c c}{\~a}o para defender oralmente qualquer etapa ou linha deste projeto. Todas as dedu{\c c}{\~o}es e dados foram validados por mim contra fontes bibliogr{\'a}ficas can{\^o}nicas.''
\end{tcolorbox}

Essa sistem{\'a}tica substitui o medo da puni{\c c}{\~a}o pela cultura da responsabilidade acad{\^e}mica, alinhando a conduta discente aos princ{\'i}pios de integridade cient{\'i}fica defendidos pelo Minist{\'e}rio da Educa{\c c}{\~a}o \cite{brasil2026referencial}.

\begin{tcolorbox}[colback=white, colframe=boxborder, arc=1mm, title={Texto pronto para o plano de ensino (copie e adapte)}]
\small
``Nesta disciplina, o uso de IA generativa segue a Escala AIAS declarada em cada enunciado. Salvo indica{\c c}{\~a}o em contr{\'a}rio, as atividades s{\~a}o N{\'i}vel AIAS \underline{\hspace{1cm}} (\underline{\hspace{4cm}}). Quando permitido, declare a ferramenta, a vers{\~a}o e a finalidade no anexo de transpar{\^e}ncia e guarde seus \emph{prompts}. Tecnologias assistivas e adapta{\c c}{\~o}es autorizadas pela institui{\c c}{\~a}o continuam permitidas e n{\~a}o se confundem com IA generativa. Em caso de d{\'u}vida, pergunte antes de entregar.''
\end{tcolorbox}
